% ============================================================
%  Section 7 — Evaluation
% ============================================================
\section{Evaluation}
\label{sec:evaluation}

The agent baseline comes from pipeline runs on the SELECTION set (30 synthetic
documents per company, 60 total) using skill version v001. Vision and latency
figures come from separate, smaller synthetic benchmarks. They describe these
fixtures and should not be interpreted as production performance.

\subsection{Metrics}

\textbf{Extractor NEM} (Normalized Exact Match): fraction of gold fields for
which the predicted value, after stripping whitespace, lowercasing, and converting
French number format to decimal, exactly matches the gold label.

\textbf{Coder Score}: binary — 1.0 only if the account code matches the gold
label and the cited string appears in a retrieved chunk or skill file. This is a
substring check; it does not establish that the citation semantically supports
the decision.

\textbf{Validator F-$\beta$} ($\beta\!=\!\sqrt{2}$): a weighted F-score where
missing an error (false negative) costs twice as much as a false alarm (false
positive). This reflects the accounting domain: an undetected arithmetic error
is worse than an unnecessary escalation.

\textbf{Vision accuracy}: fraction of key fields (supplier, document number, TTC)
correctly read from synthetic degraded images, measured on 12 generated
documents.

\subsection{LLM Health Check}

Before running any agents, a doctor script validates the model:

\begin{table}[ht]
\caption{LLM health check results (Kimi-K2.5, \texttt{reasoning\_effort=none})}
\label{tab:doctor}
\centering\scriptsize\renewcommand{\arraystretch}{1.2}
\begin{tabular}{@{}lcc@{}}
\toprule
\textbf{Check} & \textbf{Result} & \textbf{Pass?} \\
\midrule
Vision accuracy — supplier (N=12)    & 12/12 = 100\% & \checkmark \\
Vision accuracy — TTC (N=12)         & 12/12 = 100\% & \checkmark \\
Reasoning tokens in agent calls      & 0             & \checkmark \\
Extractor median latency (N=10)      & 1{,}688\,ms   & \checkmark \\
Extractor p95 latency                & 2{,}765\,ms   & \checkmark \\
JSON parse rate (text docs, N=10)    & 10/10 = 100\% & \checkmark \\
\bottomrule
\end{tabular}
\end{table}

In this sample, the vision model read the three reported fields correctly on all
12 generated images, and the text benchmark parsed JSON on all 10 documents.
These small, synthetic checks do not establish general image or extraction
accuracy. The 1.688\,s median is Extractor latency only; end-to-end latency was
not measured in this benchmark.

\subsection{Baseline Results}

\begin{table}[ht]
\caption{Baseline scores on SELECTION set, skill v001 (real runs)}
\label{tab:baseline_main}
\centering\scriptsize\renewcommand{\arraystretch}{1.25}
\begin{tabular}{@{}lrrr@{}}
\toprule
\textbf{Company} & \textbf{Coder+RAG}
                 & \textbf{Coder\,\,noRAG} & \textbf{Val.\,$F_\beta$} \\
\midrule
Company A (N=30) & 0.300 & 0.100 & 0.108 \\
Company B (N=30) & 0.205 & 0.000 & 0.000 \\
\midrule
Macro average    & 0.253 & 0.050 & 0.054 \\
\bottomrule
\end{tabular}
\end{table}

The earlier Extractor NEM values are omitted. The baseline scorer treated
unsupported or mis-mapped fields as errors; the corrected scorer is available,
but the full baseline has not been rerun, so no corrected aggregate is claimed.

\begin{figure}[ht]
\centering
\begin{tikzpicture}
\begin{axis}[
  width=\columnwidth, height=5.2cm,
  ybar, bar width=0.32cm,
  ymin=0, ymax=1.0,
  ylabel={Score},
  ylabel style={font=\scriptsize},
  symbolic x coords={Coder+RAG, Coder\,noRAG, Val.\,$F_\beta$},
  xtick=data,
  x tick label style={font=\scriptsize, rotate=12, anchor=east},
  tick label style={font=\scriptsize},
  legend style={font=\scriptsize, at={(0.97,0.97)}, anchor=north east,
                legend columns=2},
  grid=major, grid style={dashed, gray!25},
  enlarge x limits=0.15,
]
\addplot[fill=C1!70, draw=C1]   coordinates {
  (Coder+RAG,0.300)(Coder\,noRAG,0.100)(Val.\,$F_\beta$,0.108)};
\addplot[fill=C2!70, draw=C2]   coordinates {
  (Coder+RAG,0.205)(Coder\,noRAG,0.000)(Val.\,$F_\beta$,0.000)};
\legend{Company A, Company B}
\end{axis}
\end{tikzpicture}
\caption{Baseline agent scores per company on the SELECTION set (skill v001,
real pipeline runs). The with/without-RAG comparison is a small observational
ablation, not a general causal estimate.}
\label{fig:baseline_bars}
\end{figure}

\subsection{What the Numbers Mean}

\textbf{Extractor.} The previously reported 0.61 / 0.31 values are withdrawn.
The baseline scorer counted a benchmark-only reference field and failed to map
line descriptions from the production nested schema. The evaluation needs to be
rerun with the corrected scorer before an Extractor baseline can be reported.

\textbf{Coder 0.30 / 0.21 with RAG.}
These are low scores from earlier full-pipeline runs. A new paired prompt
comparison isolated Company B coding with gold extracted fields and identical
top-5 RAG context on 10 fixed SELECTION documents. Code accuracy rose from
7/10 to 9/10; the stricter correct-code-plus-supported-citation score rose
from 2/10 to 6/10. The improved prompt won four cases, lost none, and tied six.
This small selection sample is promising but is neither a full-corpus nor a
held-out TEST result. The older v001/v002 comparison remains rate-limited and
inconclusive. See \cref{sec:skillopt} and
\url{eval/results/company_b_coder_improvement.json}.

\textbf{Validator F-$\beta$ 0.11 / 0.00.}
These low scores show that validation and the baseline evaluator need further
work. We do not claim that a later code change improved them because no
comparable rerun is available. A useful next experiment should report
per-check precision and recall as well as the aggregate F-score.

The reported retrieval figures are source-level and saturated because each
company has only two or three chunks. See \cref{sec:rag} for the metric limits.
